\documentclass{article}
\usepackage[T1]{fontenc}
\usepackage{lmodern}
\usepackage{graphicx}
\usepackage{amsmath,amssymb}
\usepackage[a4paper,margin=2cm]{geometry}
\usepackage[absolute,overlay]{textpos}
\setlength{\TPHorizModule}{1mm}
\setlength{\TPVertModule}{1mm}
\usepackage[indonesian]{babel}
\usepackage{hyperref}
\setlength{\emergencystretch}{2em}
% unit: Halaman 7

\begin{document}

% === Halaman 7 (llm) ===

\begin{itemize}
    \item Supaya lebih aman, cipherteks dikelompokkan ke dalam kelompok $n$-huruf, misalnya kelompok 4-huruf:
    
    \noindent Semula: \texttt{WHPXL VDBD GL MHPEDWDQ PHUDK QDQWL PDODP}\\
    Menjadi: \texttt{WHPX LVDB DGLM HPED WDQP HUDK QDQW LPDO DP}
    
    \item Atau membuang semua spasi:
    \[ \texttt{WHPXLVDBDGLMHPEDWDQPHUDKQDQWLPDPDP} \]
    
    \item Tujuannya agar proses kriptanalisis menjadi lebih sulit dilakukan
\end{itemize}

\end{document}
